% ============================================================================
% 2. RELATED WORK — Journal version (significantly expanded)
% ============================================================================
\section{Related Work}
\label{sec:related}

\subsection{Text-to-SQL Benchmarks}

The text-to-SQL evaluation landscape has evolved rapidly.
Spider~\cite{yu2018spider} introduced cross-database evaluation with
200 databases and complex SQL queries.  BIRD~\cite{li2023bird} scaled
to 95 large databases with dirty values and external knowledge.
Spider~2.0~\cite{lei2024spider2} targets real-world enterprise workflows
with notebook-based SQL generation.  BEAVER~\cite{beaver2024} focuses
on enterprise-specific patterns.  While these benchmarks evaluate SQL
\emph{generation}, DW-Bench evaluates schema \emph{understanding}---the
structural topology reasoning that precedes and informs query construction.

\subsection{Graph Reasoning with LLMs}

Several benchmarks evaluate LLMs on graph reasoning tasks.
NLGraph~\cite{wang2023nlgraph} tests shortest path and connectivity
on synthetic graphs up to 100 nodes.
Talk-like-a-Graph~\cite{fatemi2023talk} explores different graph
serialization formats.
GraCoRe~\cite{gracore2024} benchmarks comprehension and complex reasoning
across 10 graph tasks.
GraphArena~\cite{grapharena2024} scales to million-node graphs with
computational problems like graph coloring and TSP.

DW-Bench differs from these in three key ways: (1)~it uses
\emph{heterogeneous} graphs with mixed edge types (FK + lineage),
reflecting real enterprise schemas; (2)~it includes domain-specific
subtypes (impact analysis, provenance tracing) not present in generic
graph benchmarks; and (3)~it introduces a value-level tier requiring
interactive data access.

\subsection{Tool-Augmented and Agentic LLMs}

Recent work has explored augmenting LLMs with graph-specific tools.
Think-on-Graph (ToG)~\cite{sun2024tog} interleaves LLM reasoning with
knowledge graph exploration.  Graph Chain-of-Thought~\cite{jin2024graphcot}
augments LLMs with explicit graph reasoning chains.  G-Retriever~\cite{he2024gretriever}
combines GNN-based retrieval with LLM generation for graph QA.

Our Tool-Use and ReAct-Code baselines provide purpose-built graph
algorithms (BFS, component detection, lineage traversal) as callable
tools, establishing a conservative upper bound on what tool-augmented
approaches can achieve.  We show that even with perfect tool access,
LLMs struggle with compositional multi-hop reasoning.

\subsection{Data Lineage and Provenance}

Data lineage tracking is fundamental to data governance and has been
studied extensively in the database systems community.  While lineage
systems focus on \emph{computing} provenance, DW-Bench evaluates whether
LLMs can \emph{reason about} lineage relationships expressed as graph
topology---a complementary capability relevant to LLM-assisted data
engineering workflows.

\subsection{Retrieval-Augmented Generation}

RAG~\cite{lewis2020rag} has become the standard approach for grounding
LLMs in external knowledge.  Vector-RAG uses dense retrieval
(Sentence-BERT~\cite{reimers2019sentencebert}) to inject relevant
schema fragments.  Graph-Augmented generation extends RAG by structuring
retrieved context as BFS neighborhoods from the schema graph.  DW-Bench
provides controlled comparisons across these paradigms.
